% GENERATED FILE. Edit canonical lessons, then run npm run generate:books.
% canonical-source-hash: fnv1a64:7a60878dd66a5406
% canonical-lessons: ES-C31-once-quince, ES-C31-dieciseis-diecinueve, ES-C31-teens-latinos, ES-C31-veinte

\chapter{Numbers Eleven to Twenty}
\label{ch:spanish-numbers-eleven-to-twenty}
\hlchaptermodality{40}

\begin{chapteropening}
\textbf{By the end of this chapter:} \emph{I can count from eleven to twenty in Spanish and say where the fused Latin teens stop and the ten-and-X ones begin.}

Run once to quince as fused Latin compounds, then dieciséis to diecinueve as a ten-and-X template, say why Latin's own subtractive 18 and 19 did not survive, and close on veinte.
\end{chapteropening}

\section[once — quince]{once — quince: five words with a seam you can no longer see}

\label{lesson:ES-C31-once-quince}



\begin{quote}
Spanish counts past ten in \textbf{two completely different ways}, and the join between them is a real boundary in the language, not a place someone drew a line. Today: the first five, which you simply have to learn as words.
\end{quote}

\begin{cousinweb}
\noindent
\begin{tabularx}{\linewidth}{@{}>{\raggedright\arraybackslash}X>{\raggedright\arraybackslash}X@{}}
\toprule
\textbf{Spanish} & \textbf{$\leftarrow$ Latin} \\
\midrule
\textbf{once} & \emph{undecim} (\textquotedblleft{}one-ten\textquotedblright{}) \\
\textbf{doce} & \emph{duodecim} (\textquotedblleft{}two-ten\textquotedblright{}) \\
\textbf{trece} & \emph{tredecim} (\textquotedblleft{}three-ten\textquotedblright{}) \\
\textbf{catorce} & \emph{quattuordecim} (\textquotedblleft{}four-ten\textquotedblright{}) \\
\textbf{quince} & \emph{quindecim} (\textquotedblleft{}five-ten\textquotedblright{}) \\
\bottomrule
\end{tabularx}

Every one of these \emph{was} a compound — a number plus \textbf{-decim} (\textquotedblleft{}ten,\textquotedblright{} the same root as \emph{diez}). But these five are the numbers people say most often, and heavy use wears words smooth. Centuries of everyday speech fused each pair into a single, \textbf{opaque} word: no Spanish speaker hears \textquotedblleft{}one-ten\textquotedblright{} in \emph{once} any more than an English speaker hears \textquotedblleft{}two-ten\textquotedblright{} in \emph{twelve}.

Listen for the ghost, though. The \textbf{-ce} ending on \emph{doce, trece, catorce, quince} is all that survives of Latin's \emph{-decim}. It is the same \textquotedblleft{}ten\textquotedblright{} you already say in \textbf{diez}, ground down to two letters.
\end{cousinweb}

\subsection*{Your turn}
Say these aloud:

\begin{itemize}
\raggedright
  \item \textquotedblleft{}once, doce, trece, catorce, quince\textquotedblright{}
  \item just the endings — \textquotedblleft{}-ce, -ce, -ce, -ce\textquotedblright{} — the worn-down \emph{ten}
\end{itemize}

\begin{tcolorbox}[breakable,colback=teal!4,colframe=teal!35!black,title={Before you move on}]
Count eleven to fifteen. (\emph{Once, doce, trece, catorce, quince}.) What are these five, structurally? (\textbf{Fused Latin compounds} — number plus \emph{-decim}, \textquotedblleft{}ten.\textquotedblright{}) Which piece of each word is the old \textquotedblleft{}ten\textquotedblright{}? (The \textbf{-ce} ending — \emph{diez} worn down.) Next: the numbers where the seam is still wide open.
\end{tcolorbox}
\section[dieciséis — diecinueve]{dieciséis — diecinueve: where the seam reopens}

\label{lesson:ES-C31-dieciseis-diecinueve}



\begin{quote}
You just learned five numbers whose insides had been ground smooth. Now say \textbf{fifteen} and then \textbf{sixteen} — \emph{quince}, \emph{dieciséis} — and hear the language change gear. From here you are not memorising words; you are following a rule.
\end{quote}

\begin{cousinweb}
\noindent
\begin{tabularx}{\linewidth}{@{}>{\raggedright\arraybackslash}X>{\raggedright\arraybackslash}X@{}}
\toprule
\textbf{Spanish} & \textbf{literally} \\
\midrule
\textbf{dieciséis} & \emph{diez y seis}, \textquotedblleft{}ten and six\textquotedblright{} \\
\textbf{diecisiete} & \emph{diez y siete}, \textquotedblleft{}ten and seven\textquotedblright{} \\
\textbf{dieciocho} & \emph{diez y ocho}, \textquotedblleft{}ten and eight\textquotedblright{} \\
\textbf{diecinueve} & \emph{diez y nueve}, \textquotedblleft{}ten and nine\textquotedblright{} \\
\bottomrule
\end{tabularx}

No fusion. You can hear \textbf{diez} plainly at the front of every one, and the number you already know sitting right behind it.
\end{cousinweb}

\begin{grammarlens}[title={Grammar Lens: a formula, not a list}]
This is the structural difference that makes 15/16 a real boundary in Spanish rather than a convenient place to pause:

\noindent
\begin{tabularx}{\linewidth}{@{}>{\raggedright\arraybackslash}X>{\raggedright\arraybackslash}X>{\raggedright\arraybackslash}X@{}}
\toprule
\textbf{} & \textbf{11-15} & \textbf{16-19} \\
\midrule
how you get the word & \textbf{remember} it & \textbf{build} it \\
the \textquotedblleft{}ten\textquotedblright{} inside & worn to \emph{-ce} & still spelled \emph{dieci-} \\
can you hear the parts? & no & yes \\
\bottomrule
\end{tabularx}

So \emph{once} through \emph{quince} are \textbf{vocabulary}, and \emph{dieciséis} through \emph{diecinueve} are \textbf{grammar}: one template, \emph{dieci- + digit}, filled four times. Learn the template once and all four are already yours.
\end{grammarlens}

\subsection*{Your turn}
Say these aloud:

\begin{itemize}
\raggedright
  \item \textquotedblleft{}dieciséis, diecisiete, dieciocho, diecinueve\textquotedblright{}
  \item the seam out loud — \textquotedblleft{}diez y seis\textquotedblright{} then \textquotedblleft{}dieciséis\textquotedblright{}
  \item the boundary — \textquotedblleft{}quince\textquotedblright{} (remembered), \textquotedblleft{}dieciséis\textquotedblright{} (built)
\end{itemize}

\begin{tcolorbox}[breakable,colback=teal!4,colframe=teal!35!black,title={Before you move on}]
Count sixteen to nineteen. (\emph{Dieciséis, diecisiete, dieciocho, diecinueve}.) Are these fused or transparent? (\textbf{Transparent} — plain \textquotedblleft{}ten and X.\textquotedblright{}) Which of the two groups do you have to memorise, and which can you generate? (\textbf{Memorise 11-15}; \textbf{generate 16-19} from \emph{dieci-} plus the digit.)
\end{tcolorbox}
\section[duodēvīgintī, ūndēvīgintī]{duodēvīgintī — the quirk Spanish threw away}

\label{lesson:ES-C31-teens-latinos}



\begin{quote}
\emph{Dieciocho} and \emph{diecinueve} look like the least interesting words in the chapter — just two more ten-and-X compounds. They are actually the most interesting, because of what Latin said instead.
\end{quote}

\subsection*{What to know first}
\begin{itemize}
\raggedright
  \item dieciséis — diecinueve — the \emph{dieci-} plus digit template.
\end{itemize}

\begin{culture}
You would expect \emph{dieciocho} and \emph{diecinueve} to carry on whatever Latin did for 18 and 19. Latin did something else entirely. Its everyday words for those two numbers were \textbf{subtractive} — counted \emph{down} from twenty:

\noindent
\begin{tabularx}{\linewidth}{@{}>{\raggedright\arraybackslash}X>{\raggedright\arraybackslash}X>{\raggedright\arraybackslash}X@{}}
\toprule
\textbf{number} & \textbf{Latin} & \textbf{literally} \\
\midrule
18 & \textbf{duodēvīgintī} & \textquotedblleft{}\textbf{two from} twenty\textquotedblright{} \\
19 & \textbf{ūndēvīgintī} & \textquotedblleft{}\textbf{one from} twenty\textquotedblright{} \\
\bottomrule
\end{tabularx}

(You still meet this habit in English clock-talk: \textquotedblleft{}ten \textbf{to} four.\textquotedblright{})

Spanish did not inherit either word. It built \emph{dieciocho} and \emph{diecinueve} fresh, on exactly the same simple template as 16 and 17 — quietly \textbf{regularising away} an irregularity that existed in its own parent language. Languages are usually described as wearing down and getting messier; here is a clean case of one tidying up.
\end{culture}

\subsection*{Your turn}
Say these aloud:

\begin{itemize}
\raggedright
  \item Latin's way — \textquotedblleft{}two from twenty\textquotedblright{} — then Spanish's — \textquotedblleft{}dieciocho\textquotedblright{}
  \item \textquotedblleft{}dieciocho, diecinueve\textquotedblright{} — addition, not subtraction
\end{itemize}

\begin{tcolorbox}[breakable,colback=teal!4,colframe=teal!35!black,title={Before you move on}]
How did Latin say 18 and 19? (\textbf{Subtractively} — \emph{duodēvīgintī}, \textquotedblleft{}two from twenty,\textquotedblright{} and \emph{ūndēvīgintī}, \textquotedblleft{}one from twenty.\textquotedblright{}) Did Spanish keep that? (\textbf{No} — it built \emph{dieciocho} and \emph{diecinueve} on the regular ten-and-X template.) So which language is the irregular one here? (\textbf{Latin} — Spanish simplified.)
\end{tcolorbox}
\section[veinte]{veinte — the number the teens were counting toward}

\label{lesson:ES-C31-veinte}



\begin{quote}
Latin's own 18 and 19 were counted \emph{down from twenty}. So it is worth meeting twenty itself — the number the whole run has been aiming at.
\end{quote}

\begin{cousinweb}
\textbf{veinte} (\textquotedblleft{}twenty\textquotedblright{}) $\leftarrow$ Latin \textbf{vīgintī}. It was worn down by ordinary sound change — the \emph{g} softened away, the vowels shifted — but notice what did \textbf{not} happen to it: it was never \textbf{fused} out of two separate words the way \emph{once} through \emph{quince} were. \emph{Vīgintī} was already a single word for a single round number, so there is no seam inside \emph{veinte} to look for.

Three different histories, then, in one short chapter:

\noindent
\begin{tabularx}{\linewidth}{@{}>{\raggedright\arraybackslash}X>{\raggedright\arraybackslash}X@{}}
\toprule
\textbf{words} & \textbf{what happened} \\
\midrule
\emph{once - quince} & two words \textbf{fused} into one \\
\emph{dieciséis - diecinueve} & two words \textbf{built} into one, seam visible \\
\emph{veinte} & one word, simply \textbf{worn} \\
\bottomrule
\end{tabularx}
\end{cousinweb}

\subsection*{Your turn}
Say these aloud:

\begin{itemize}
\raggedright
  \item \textquotedblleft{}once, doce, trece, catorce, quince\textquotedblright{} — the five fused Latins
  \item \textquotedblleft{}dieciséis, diecisiete, dieciocho, diecinueve\textquotedblright{} — built, not fused
  \item \textquotedblleft{}veinte\textquotedblright{} — twenty, worn but never fused
\end{itemize}

\begin{tcolorbox}[breakable,colback=teal!4,colframe=teal!35!black,title={Before you move on}]
Count eleven to twenty. Where is the boundary between remembering and building, and why is it there? (\textbf{After \emph{quince}} — 11-15 are fused Latin compounds, 16-19 are \emph{dieci-} plus a digit.) Did Spanish inherit Latin's subtractive 18 and 19? (\textbf{No} — it regularised them.) Was \emph{veinte} fused from two words? (\textbf{No} — Latin \emph{vīgintī} was already one word; it was only worn down.)
\end{tcolorbox}
